\subsection{Proof of Theorem 1: first part}

The implication (i) \(\Rightarrow\) (ii) has already been proved (X, p.~157, prop. 5). The implication (ii) \(\Rightarrow\) (iii) is evident; likewise for (iv) \(\Rightarrow\) (v), since \(\beta_{\mathbf{M}}^{x}\) is identified with \(\alpha_{\mathbf{M}}^{x} \otimes 1_{\mathbf{A}/\mathfrak{x}}\) .

To show that (iv) implies (iii), consider the homomorphism \((\alpha_{M}^{x})_{1}\) induced on the degree-1 components. Let E denote the graded A-module \(A[X_{1},\ldots,X_{n}]\) ; the A-module \(E_{1}\) is free with basis \(X_{1},\ldots,X_{n}\) and \(\mathfrak{d}_{1}\) is the sub-A-module of \(E_{1}\) generated by the elements \((x_{i}X_{j}-x_{j}X_{i})\) for \(1\leqslant i<j\leqslant n\) . Consequently, \(((E/\mathfrak{d})\otimes_{A}M)_{1}\) is identified with \(K_{1}(x,M)/B_{1}(K.(x,M))\) , the homomorphism \((\alpha_{M}^{x})_{1}\) being identified with the map from \(K_{1}(x,M)/B_{1}(K.(x,M))\) over \(B_{0}(K.(x,M))\) induced by \(d_{1}\) . Thus the vanishing of \(H_{1}(x,M)\) is equivalent to the injectivity of \((\alpha_{M}^{x})_{1}\) , whence the implication (iv) \(\Rightarrow\) (iii).

To prove that (iii) implies (iv), we shall use the following lemma:

Lemma 4. --- Let \(\mathbf{A}_{0}\) be the ring \(\mathbf{Z}[\mathrm{T}_{1},\ldots,\mathrm{T}_{n}]\) , and let \(u:\mathbf{A}_{0}[\mathrm{X}_{1},\ldots,\mathrm{X}_{n}] \to \mathbf{A}_{0}[\mathrm{U}]\) be the homomorphism of \(\mathbf{A}_{0}\) -algebras such that \(u(\mathbf{X}_{i}) = \mathbf{T}_{i}\) U. The kernel of \(u\) is the ideal \(\mathfrak{d}_{0}\) of \(\mathbf{A}_{0}[\mathbf{X}_{1},\ldots,\mathbf{X}_{n}]\) generated by the elements ( \(\mathrm{T}_{i}\mathbf{X}_{j}-\mathrm{T}_{j}\mathbf{X}_{i}\) ) for \(1 \leqslant i < j \leqslant n\) . If t is the ideal of \(\mathbf{A}_{0}\) generated by ( \(\mathrm{T}_{1},\ldots,\mathrm{T}_{n}\) ), \(u\) induces an isomorphism

\[
\overline {{u}}: \mathrm{A} _ {0} [ \mathrm{X} _ {1}, \dots , \mathrm{X} _ {n} ] / \mathfrak {d} _ {0} \rightarrow \bigoplus_ {r \geqslant 0} t ^ {r}.
\]

It obviously suffices to prove the first assertion. For every sequence of natural integers \(\alpha = (\alpha_{1},\ldots,\alpha_{n})\) and every integer \(k\in[0,n]\) , denote by \(\mathbf{P}_{\alpha,k}\) the monomial

\[
\mathrm{T} _ {1} ^ {\alpha_ {1}} \dots \mathrm{T} _ {k} ^ {\alpha_ {k}} \mathrm{X} _ {k + 1} ^ {\alpha_ {k + 1}} \dots \mathrm{X} _ {n} ^ {\alpha_ {n}};
\]

let N be the sub-Z-module of \(A_{0}[X_{1}, \ldots, X_{n}]\) generated by the \(P_{\alpha,k}\) for \(\alpha \in N^{n}\) and \(0 \leqslant k \leqslant n\) . We shall show that the restriction of u to N is injective and that \(A_{0}[X_{1}, \ldots, X_{n}] = \mathfrak{d}_{0} + N\) ; since one has \(d_{0} \subset Ker u\) , this will imply the lemma.

Observe that N is generated by the set S consisting of \(P_{0,0} = 1\) and of those \(P_{\alpha,k}\) for which \(\alpha_k \neq 0\) . To prove the injectivity of the restriction of u to N, it suffices to show that two distinct elements of S have as images under u distinct monomials in \(A_0[U]\) . Now one has \(u(P_{\alpha,k}) = T^\alpha U^{\sum_{i \geq k} \alpha_i}\) , so that the equality \(u(P_{\alpha,k}) = u(P_{\alpha',k'})\) implies \(\alpha = \alpha'\) and \(\sum_{i \geq k} \alpha_i = \sum_{i \geq k'} \alpha_i\) . Suppose that \(P_{\alpha,k}\) and \(P_{\alpha',k'}\) belong to S. If \(\alpha = 0\) , one then has \(k = k' = 0\) ; if \(\alpha \neq 0\) , we obtain \(k = k'\) since \(\alpha_k\) and \(\alpha_{k'}\) are nonzero, hence the result.

Let us show that every monomial \(\mathbf{T}^{\alpha}\mathbf{X}^{\beta}\in\mathbf{A}_{0}[\mathbf{X}_{1},..., \mathbf{X}_{n}]\) is congruent modulo \(\mathfrak{d}_{0}\) to a \(\mathbf{P}_{\eta,k}\) . For every sequence \(\lambda\in\mathbf{N}^{n}\) , we denote by \(i(\lambda)\) (resp. \(j(\lambda)\) ) the smallest (resp. the largest) integer \(k\in[1,n]\) such that \(\lambda_{k}\neq0\) . In the set of monomials \(\mathbf{T}^{\gamma}\mathbf{X}^{\delta}\) that are congruent to \(\mathbf{T}^{\alpha}\mathbf{X}^{\beta}\) modulo \(\mathfrak{d}_{0}\) let us choose one for which the rational integer \(j(\gamma)-i(\delta)\) is minimal; let us show that we then have \(j(\gamma)-i(\delta)<0\) . Suppose that we have \(j(\gamma)\geqslant i(\delta)\) ; denote by \(j=j(\gamma)\) , \(i=i(\delta)\) , and let \(\varepsilon=\inf(\gamma_{j},\delta_{i})\) . Since \((\mathbf{T}_{j}^{\varepsilon}\mathbf{X}_{i}^{\varepsilon}-\mathbf{T}_{i}^{\varepsilon}\mathbf{X}_{j}^{\varepsilon})\) is divisible by \((\mathbf{T}_{j}\mathbf{X}_{i}-\mathbf{T}_{i}\mathbf{X}_{j})\) , hence belongs to \(\mathfrak{d}_{0}\) , we see that \(\mathbf{T}^{\gamma}\mathbf{X}^{\delta}\) is congruent modulo \(\mathfrak{d}_{0}\) to \(\mathbf{T}^{\gamma'}\mathbf{X}^{\delta'}\) , where \(\gamma_{i}'=\gamma_{i}+\varepsilon\) , \(\gamma_{j}'=\gamma_{j}-\varepsilon\) , \(\gamma_{k}'=\gamma_{k}\) for \(k\neq i,j\) , and \(\delta_{i}'=\delta_{i}-\varepsilon\) , \(\delta_{j}'=\delta_{j}+\varepsilon\) , \(\delta_{k}'=\delta_{k}\) for \(k\neq i,j\) . Since \(\gamma_{j}'\) or \(\delta_{i}'\) is zero, we have \(j(\gamma')-i(\delta')<j(\gamma)-i(\delta)\) which contradicts the minimal character of \(j(\gamma)-i(\delta)\) .

Consequently we have \(j(\gamma) < i(\delta)\) , whence \(\mathbf{T}^{\gamma}\mathbf{X}^{\delta} \in \mathbf{N}\) , which completes the proof of the lemma.

Let us prove that (iii) implies (iv). Consider the ring \(\mathbf{A}_0 = \mathbf{Z}[\mathrm{T}_1, \ldots, \mathrm{T}_n]\) and the ideal t of \(\mathbf{A}_0\) generated by \(\mathrm{T}_1, \ldots, \mathrm{T}_n\) . Endow M with the structure of \(\mathbf{A}_0\) -module for which \(\mathrm{T}_i m = x_i m\) for \(m \in \mathbf{M}\) , \(1 \leqslant i \leqslant n\) . By X, p.~155, \(\mathrm{H}_i(x, \mathbf{M})\) is canonically identified with \(\mathbf{H}_i(\mathbf{T}, \mathbf{M})\) .

Since the sequence \(\mathbf{T}\) is regular for \(\mathbf{A}_{0}\) it follows from implication (i) \(\Rightarrow\) (ii) that the complex \(\mathbf{K}.\left(\mathbf{T},\mathbf{A}_{0}\right)\) defines a free resolution of the \(\mathbf{A}_{0}\) -module \(\mathbf{A}_{0}/t\) . Note that this one identifies with \(\mathbf{Z}\) , equipped with the structure of \(\mathbf{A}_{0}\) -module such that \(\mathrm{T}_{i}\mathbf{Z}=0\) for \(1\leqslant i\leqslant n\) . Thus condition (iii) is equivalent to \(\operatorname{Tor}_{1}^{\mathbf{A}_{0}}(\mathbf{M},\mathbf{Z})=0\) .

Let us show that (iii) implies \(\mathrm{Tor}_{1}^{\mathbf{A}_{0}}(\mathbf{M},\mathbf{A}_{0}/t^{r})=0\) for all \(r\geqslant1\) . This follows from the preceding for \(r=1\) . In the general case, consider the exact sequence

\[
0 \rightarrow t ^ {r} / t ^ {r + 1} \rightarrow A _ {0} / t ^ {r + 1} \rightarrow A _ {0} / t ^ {r} \rightarrow 0.\tag{16}
\]

The \(A_{0}\) -module \(t^{r}/t^{r+1}\) is isomorphic to a finite product of copies of Z; we therefore have \(\operatorname{Tor}_{1}^{\mathbf{A}_{0}}(\mathbf{M}, t^{r}/t^{r+1}) = 0\) . From this, by induction on r, we deduce that

\[
\operatorname{Tor} _ {1} ^ {\mathrm{A} _ {0}} \left(\mathrm{M}, \mathrm{A} _ {0} / t ^ {r}\right) = 0 \quad \text { for   all } r.
\]

The exact sequence (16) therefore yields, by tensor product with M, an exact sequence

\[
0 \rightarrow (t ^ {r} / t ^ {r + 1}) \otimes_ {\mathrm{A} _ {0}} \mathrm{M} \rightarrow \mathrm{M} / x ^ {r + 1} \mathrm{M} \rightarrow \mathrm{M} / x ^ {r} \mathrm{M} \rightarrow 0\tag{17}
\]

whence an isomorphism of \((\mathfrak{t}^r /\mathfrak{t}^{r + 1})\otimes_{\mathbf{A}_0}\mathbf{M}\) over \(\mathfrak{x}^r\mathbf{M} / \mathfrak{x}^{r + 1}\mathbf{M}\) . Considering the exact sequence \(0\to \mathfrak{t}^{r + 1}\to \mathfrak{t}^r\to \mathfrak{t}^r /\mathfrak{t}^{r + 1}\to 0\) , one then sees by induction on \(r\) that the map \(m_r:\mathfrak{t}^r\otimes_{\mathbf{A}_0}\mathbf{M}\rightarrow \mathfrak{x}^r\mathbf{M}\) , induced by the action of \(\mathbf{A}_0\) to \(\mathbf{M}\) is an isomorphism.

To prove (iv), it remains to observe that the diagram

\[
\begin{array}{c} \left(\mathrm{A} _ {0} \left[ \mathrm{X} _ {1}, \dots , \mathrm{X} _ {n} \right] / \mathfrak {d} _ {0}\right) \otimes_ {\mathrm{A} _ {0}} \mathrm{M} \xrightarrow {\bar {u} \otimes 1 _ {\mathrm{M}}} \bigoplus_ {r \geqslant 0} \left(\mathrm{t} ^ {r} \otimes_ {\mathrm{A} _ {0}} \mathrm{M}\right) \\ \Bigg \downarrow^ {e} \\ ^ {\prime} (\mathrm{A} [ \mathrm{X} _ {1}, \dots , \mathrm{X} _ {n} ] / \mathfrak {d}) \otimes_ {\mathrm{A}} \mathrm{M} \xrightarrow {\alpha_ {\mathrm{M}} ^ {*}} \bigoplus_ {r \geqslant 0} x ^ {r} \mathrm{M} \end{array}
\]

where e is the canonical isomorphism of extension of scalars (II, p.~85, prop. 6), is commutative, and to apply Lemma 4.

